% !TeX program = xelatex
% SPDX-License-Identifier: MIT
% Copyright (c) 2026 Ziyu Peng and 刘毅涵
% Author: https://github.com/pzy54154631-hub?tab=repositories
\documentclass[UTF8,a4paper,fontset=fandol]{ctexart}
\usepackage[margin=2.5cm]{geometry}
\usepackage[hidelinks]{hyperref}
\setcounter{tocdepth}{2}
\setcounter{secnumdepth}{2}

\title{课程报告的结构练习}
\author{Ziyu Peng \and 刘毅涵}
\date{}

\begin{document}
\maketitle
\tableofcontents
\clearpage

\section{研究问题}
\label{sec:question}
本例只演示文档结构，不包含实际研究结论。
材料范围见第~\ref{sec:materials}~节。

\section{材料与方法}
\label{sec:methods}
\subsection{材料范围}
\label{sec:materials}
这里说明准备使用哪些材料，以及为什么选择它们。

这里另起一段，说明材料的限制。
\footnote{本脚注仅用于演示补充说明。}

\subsection{整理步骤}
\begin{enumerate}
  \item 记录材料出处。
  \item 按主题进行整理。
  \item 检查说明与材料是否一致。
\end{enumerate}

\section{小结}
问题在第~\pageref{sec:question}~页提出。
当前文档已经可以自动更新编号。

\section*{致谢}
\phantomsection
\addcontentsline{toc}{section}{致谢}
这里放置不编号的致谢内容。
\end{document}
